\subsection{Lattice ordered groups}

Recall that an ordered set in which every non empty finite subset has a supremum and an infimum is called a lattice (E, III, p.~13). It is clear that a product of lattice ordered groups, and in particular a product of totally ordered groups, is a lattice ordered group. In contrast a subgroup of a lattice ordered group is not necessarily lattice ordered.

Thus, in the product ordered group \(\mathbf{Z} \times \mathbf{Z}\) , the «antidiagonal» (set of pairs \((n, n')\) such that \(n + n' = 0\) ) is ordered with the discrete ordering, and hence is not a lattice ordered group. * The additive group of polynomials in one real variable (VI, p.~1, example 2) is a filtered group (since both \(p(x)\) and \(q(x)\) are less than \((p(x))^2 + (q(x))^2 + 1\) ) which one can show is not lattice ordered. *

PROPOSITION 7. --- If \(x\) and \(y\) are two elements of an ordered group \(G\) , and if one of the elements \(\inf(x, y)\) , \(\sup(x, y)\) exists, then so does the other, and \(x + y = \inf(x, y) + \sup(x, y)\) .

Indeed, according to the relations (1) and (2) (VI, p.~9) we have

\[
\sup (a - x, a - y) = a + \sup (- x, - y) = a - \inf (x, y),
\]

and it suffices to take \(a = x + y\) .

PROPOSITION 7 (DIV). --- If \(a, b \in \mathbf{K}^{*}\) , and if \(d\) is a gcd of \(a\) and \(b\) , and \(m\) an lcm of \(a\) and \(b\) , then the product dm is an associate of \(ab\) .

PROPOSITION 8. --- Let P be the set of positive elements of an ordered group G. For G to be lattice ordered, it is necessary and sufficient that G = P - P, and that in addition P, equipped with the induced ordering, satisfies one or other of the following conditions :

\begin{enumerate}
\def\labelenumi{\alph{enumi})}
\item
  Each pair of elements of \(\mathbf{P}\) has a supremum in \(\mathbf{P}\) .
\item
  Each pair of elements of \(\mathbf{P}\) has an infimum in \(\mathbf{P}\) .
\end{enumerate}

The necessity of these conditions is obvious: indeed the relation \(G = P - P\) says that G is filtered (VI, p.~4, Prop. 4); on the other hand the supremum and infimum in G of two elements of P are positive, so are also their supremum and infimum in P.

Conversely, note first that under hypothesis a) (resp. b)), every pair of elements x, y of P has a supremum (resp. infimum) in G equal to its supremum a (resp. infimum b) in P. This is obvious for a, since any upper bound for x and y is positive; for b, let \(z \in G\) be a lower bound for x and y; then there exists \(u \in P\) such that \(z + u \in P\) , since G = P - P; now \(\inf_{P}(x + u, y + u)\) is greater than \(b + u\) , and so is of the form \(b + c + u\) ( \(c \geqslant 0\) ); since \(b + c\) is less than x and less than y, we have c = 0; hence \(\inf_{P}(x + u, y + u) = b + u\) , which implies \(z + u \leqslant b + u\) , thus \(z \leqslant b\) and b is certainly the infimum of x and y in G. Now if x and y are arbitrary elements of G we translate them into P: let \(v \in P\) be such that \(x + v\) and \(y + v\) are positive (VI, p.~5, Prop. 5); under hypothesis a) (resp. b))

there exists a supremum (resp. infimum) for \(x + v\) and \(y + v\) in P, and hence also in G by what we have just seen; by translation x and y have a supremum (resp. infimum) in G; the existence of one of these implies the other, by Prop. 7, and this shows that the conditions are sufficient.

